\begin{theorem}[Finite-rank kernelization of off-critical-line zero-scan profiles: from boundary scalars to \texorpdfstring{$\kappa$}{kappa}-dimensional reproducing kernel Hilbert spaces]\label{thm:xi-basepoint-scan-finite-rank-rkhs}
Following the notation of Definition~\ref{def:xi-basepoint-scan-profile} and Proposition~\ref{prop:xi-basepoint-depth-closed-form}, let
\[
a_k:=1+\delta_k,\qquad
p_k:=\gamma_k-\mathrm{i}a_k\in\CC_{-},\qquad
w_k:=4\delta_k\qquad(1\le k\le \kappa).
\]
Define the feature map \(v:\RR\to\CC^{\kappa}\) by
\[
v(x)_k:=\frac{\sqrt{w_k}}{x-p_k},
\]
and define the kernel
\[
K(x,y):=\langle v(x),v(y)\rangle_{\CC^{\kappa}}
=\sum_{k=1}^{\kappa}\frac{w_k}{(x-p_k)(y-\overline{p_k})}
=\sum_{k=1}^{\kappa}\frac{w_k}{(x-\gamma_k+\mathrm{i}a_k)(y-\gamma_k-\mathrm{i}a_k)}.
\]
Then \(K\) is a positive definite kernel, and its diagonal strictly equals the scan profile:
\[
\boxed{\;K(x,x)=H(x)\qquad(x\in\RR).\;}
\]
Consequently, the scan profile \(H\) can be precisely kernelized as the diagonal of a finite-rank kernel.
Furthermore, \(\mathcal{H}:=\overline{\Span}\{v(x):x\in\RR\}\subseteq\CC^\kappa\) is a \(\kappa\)-dimensional Hilbert space under the natural inner product (equivalently: \(K\) is the reproducing kernel of a \(\kappa\)-dimensional RKHS generated by one-parameter data).
\end{theorem}
\begin{proof}
By definition \(K(x,y)=\langle v(x),v(y)\rangle\), for any finite point set \(\{x_i\}_{i=1}^{n}\subset\RR\) and coefficients \(\{c_i\}_{i=1}^{n}\subset\CC\) we have
\[
\sum_{i,j=1}^{n}c_i\overline{c_j}K(x_i,x_j)
=\left\|\sum_{i=1}^{n}c_i v(x_i)\right\|_{\CC^\kappa}^2\ \ge\ 0,
\]
hence \(K\) is positive definite.
The diagonal identity follows directly from \(|x-p_k|^2=(x-\gamma_k)^2+a_k^2\):
\[
K(x,x)=\sum_{k=1}^{\kappa}\frac{w_k}{|x-p_k|^2}
=\sum_{k=1}^{\kappa}\frac{4\delta_k}{(x-\gamma_k)^2+(1+\delta_k)^2}
=H(x).
\]
Finally, to prove \(\dim\mathcal{H}=\kappa\), take any \(\kappa\) distinct real points \(x_1,\dots,x_\kappa\).
The matrix \(\bigl(\frac{1}{x_i-p_k}\bigr)_{1\le i,k\le\kappa}\) is a Cauchy matrix, whose determinant is strictly nonzero by the Cauchy identity (since \(x_i\in\RR\), \(p_k\in\CC_{-}\) are mutually distinct), hence \(\{v(x_i)\}_{i=1}^{\kappa}\) are linearly independent and span \(\CC^\kappa\).
Thus \(\mathcal{H}=\CC^\kappa\).
Viewing \(K\) as the inner-product kernel of the feature map \(v\) yields the corresponding \(\kappa\)-dimensional RKHS description. This completes the proof.
\end{proof}

\begin{proposition}[Kernel uniqueness of scan profile and Cauchy pole skeleton]\label{prop:xi-basepoint-scan-kernel-uniqueness}
Following the notation of Theorem~\ref{thm:xi-basepoint-scan-finite-rank-rkhs}, let the scan profile be
\(
H(x)=\sum_{k=1}^{\kappa}\frac{w_k}{|x-p_k|^2}
\).
Given a matrix \(C=(c_{k\ell})\in\CC^{\kappa\times\kappa}\), define the kernel
\[
\widetilde K(x,y):=\sum_{k,\ell=1}^{\kappa}\frac{c_{k\ell}}{(x-p_\ell)(y-\overline{p_k})}\qquad(x,y\in\RR).
\]
If diagonal consistency \(\widetilde K(x,x)=H(x)\) holds for all \(x\in\RR\), then necessarily
\[
C=\mathrm{diag}(w_1,\dots,w_\kappa),
\qquad
\widetilde K(x,y)=\sum_{k=1}^{\kappa}\frac{w_k}{(x-p_k)(y-\overline{p_k})}=K(x,y).
\]
Furthermore, if \(C\) is Hermitian positive semidefinite, then \(\widetilde K\) is a positive definite kernel; when \(w_k>0\), this kernel has rank \(\kappa\).
\end{proposition}
\begin{proof}
Let \(W:=\mathrm{diag}(w_1,\dots,w_\kappa)\) and set \(D:=C-W\).
By diagonal consistency, we obtain a rational identity for all \(x\in\RR\)
\[
0=\widetilde K(x,x)-K(x,x)
=\sum_{k,\ell=1}^{\kappa}\frac{D_{k\ell}}{(x-p_\ell)(x-\overline{p_k})}.
\]
The right-hand side is a rational function in the complex variable \(z\); since it vanishes identically on the real axis, it vanishes identically on \(\CC\).
Taking the residue of this rational function at \(z=p_j\) (\(1\le j\le \kappa\)) yields
\[
0=\operatorname*{Res}_{z=p_j}\sum_{k,\ell=1}^{\kappa}\frac{D_{k\ell}}{(z-p_\ell)(z-\overline{p_k})}
=\sum_{k=1}^{\kappa}\frac{D_{kj}}{p_j-\overline{p_k}}.
\]
Let \(M\in\CC^{\kappa\times\kappa}\) be the Cauchy matrix \(M_{kj}:=\frac{1}{p_j-\overline{p_k}}\).
Since \(\{p_j\}\subset\CC_{-}\) and \(\{\overline{p_k}\}\subset\CC_{+}\) are distinct, the Cauchy determinant identity gives \(\det M\ne 0\), hence each column vector \((D_{kj})_{k}\) must be zero, forcing \(D=0\), i.e., \(C=W\).
\par
The final statement follows from the standard Gram representation: if \(C\succeq 0\), then \(\widetilde K(x,y)=\langle C^{1/2}u(x),C^{1/2}u(y)\rangle\) (where \(u(x)_k=(x-p_k)^{-1}\)) is a positive definite kernel.
When \(w_k>0\), \(C=W\) is positive definite; combined with the non-degeneracy of the Cauchy matrix in the proof of Theorem~\ref{thm:xi-basepoint-scan-finite-rank-rkhs}, the rank equals \(\kappa\). This completes the proof.
\end{proof}

\begin{proposition}[Cauchy--Vandermonde closed form of anchoring determinant: fully explicit factorization under full-rank anchor sets]\label{prop:xi-basepoint-scan-anchor-det-cauchy-vandermonde}
Under the setup of Theorem~\ref{thm:xi-basepoint-scan-finite-rank-rkhs}, take a distinct real anchor set
\(
A=\{x_1,\dots,x_n\}\subset\RR
\)
and let \(V\in\CC^{n\times\kappa}\) be
\[
V_{ik}:=v(x_i)_k=\frac{\sqrt{w_k}}{x_i-p_k}.
\]
Denote its Gram matrix
\[
G_A:=(K(x_i,x_j))_{1\le i,j\le n}=VV^\ast.
\]
Then for any \(n\le \kappa\) we have the strict Cauchy--Binet expansion
\[
\boxed{\;
\det G_A
=\sum_{\substack{I\subseteq\{1,\dots,\kappa\}\\ |I|=n}}
\Bigl(\prod_{k\in I}w_k\Bigr)\cdot
\left|\det\!\left(\frac{1}{x_i-p_k}\right)_{\substack{1\le i\le n\\ k\in I}}\right|^2\; }.
\]
In particular, when \(n=\kappa\) (full-rank anchoring),
\[
\det G_A=\Bigl(\prod_{k=1}^{\kappa}w_k\Bigr)\cdot
\left|\det\!\left(\frac{1}{x_i-p_k}\right)_{1\le i,k\le\kappa}\right|^2
=4^\kappa\Bigl(\prod_{k=1}^{\kappa}\delta_k\Bigr)\,
\frac{|\Delta(X)|^2\,|\Delta(P)|^2}{\prod_{i,k}|x_i-p_k|^2},
\]
where \(\Delta(X):=\prod_{1\le i<j\le \kappa}(x_j-x_i)\), \(\Delta(P):=\prod_{1\le i<j\le \kappa}(p_j-p_i)\).
\end{proposition}
\begin{proof}
By \(G_A=VV^\ast\), applying Cauchy--Binet for \(n\le\kappa\) gives
\[
\det(VV^\ast)
=\sum_{|I|=n}\bigl|\det(V_{(:,I)})\bigr|^2.
\]
Moreover, \(V_{(:,I)}=\bigl(\frac{1}{x_i-p_k}\bigr)_{k\in I}\cdot \mathrm{diag}(\sqrt{w_k})_{k\in I}\), so
\(
|\det(V_{(:,I)})|^2
=(\prod_{k\in I}w_k)\cdot\bigl|\det(\frac{1}{x_i-p_k})\bigr|^2
\), yielding the first formula.
\par
When \(n=\kappa\), the Cauchy determinant identity gives
\[
\det\!\left(\frac{1}{x_i-p_k}\right)_{1\le i,k\le\kappa}
=\frac{\Delta(X)\Delta(P)}{\prod_{i,k}(x_i-p_k)}.
\]
Taking absolute value squared and substituting \(w_k=4\delta_k\) yields the conclusion. This completes the proof.
\end{proof}

\begin{corollary}[Resultant expression of full-rank anchoring determinant: denominator equals \texorpdfstring{$\prod_k|y(p_k)|^2$}{prod|y(pk)|^2}]\label{cor:xi-basepoint-scan-anchor-det-resultant}
In the full-rank case \(n=\kappa\) of Proposition~\ref{prop:xi-basepoint-scan-anchor-det-cauchy-vandermonde}, let the anchor polynomial be
\[
y(z):=\prod_{i=1}^{\kappa}(z-x_i)
\qquad(\text{monic with distinct roots}).
\]
Let the pole polynomial be
\[
\mathcal{A}(z):=\prod_{k=1}^{\kappa}(z-p_k)(z-\overline{p_k})\in\RR[z].
\]
Define the resultant
\(
\operatorname{Res}(y,\mathcal{A}):=\prod_{\mathcal{A}(\zeta)=0}y(\zeta)
\)
(counted with multiplicity).
Then we have the strict identity
\[
\boxed{\;
\operatorname{Res}(y,\mathcal{A})
=\prod_{k=1}^{\kappa}y(p_k)y(\overline{p_k})
=\prod_{k=1}^{\kappa}|y(p_k)|^{2}\ >\ 0\;}
\]
and moreover
\[
\boxed{\;
\det G_A
=\Bigl(\prod_{k=1}^{\kappa}w_k\Bigr)\,
\frac{|\Delta(X)|^2\,|\Delta(P)|^2}{\operatorname{Res}(y,\mathcal{A})}\; }.
\]
\end{corollary}
\begin{proof}
By definition,
\(
\operatorname{Res}(y,\mathcal{A})
=\prod_{k=1}^{\kappa}y(p_k)y(\overline{p_k})
\).
Since \(y\in\RR[z]\), we have \(y(\overline{p_k})=\overline{y(p_k)}\), so the product equals \(\prod_k|y(p_k)|^2\) and is strictly positive (because \(p_k\notin\RR\) ensures \(y(p_k)\ne 0\)).
\par
On the other hand, by the full-rank closed form of Proposition~\ref{prop:xi-basepoint-scan-anchor-det-cauchy-vandermonde},
\[
\det G_A=\Bigl(\prod_{k=1}^{\kappa}w_k\Bigr)\,
\frac{|\Delta(X)|^2\,|\Delta(P)|^2}{\prod_{i,k}|x_i-p_k|^2}.
\]
Note that
\(
\prod_{i,k}|x_i-p_k|^2
=\prod_{k=1}^{\kappa}\prod_{i=1}^{\kappa}|x_i-p_k|^2
=\prod_{k=1}^{\kappa}|y(p_k)|^2
=\operatorname{Res}(y,\mathcal{A})
\); substituting this back completes the proof.
\end{proof}

\begin{proposition}[Principal-angle decomposition: exact geometric meaning of anchoring mutual information and sine representation of two-point separation]\label{prop:xi-basepoint-scan-anchor-principal-angles}
Following the setup of Proposition~\ref{prop:xi-basepoint-scan-anchor-det-cauchy-vandermonde}, for \(x_i\in A\) let
\[
u_i:=\frac{v(x_i)}{\|v(x_i)\|},\qquad
\|v(x_i)\|^2=K(x_i,x_i)=H(x_i),
\]
and let the correlation matrix be \(R_A:=(\langle u_i,u_j\rangle)_{1\le i,j\le n}\).
Then
\[
\boxed{\;
\det G_A=\Bigl(\prod_{i=1}^{n}H(x_i)\Bigr)\det R_A\; }.
\]
If we perform Gram--Schmidt orthogonalization on \(\{u_i\}\) and denote \(\theta_j\in(0,\pi/2]\) as the principal angle between \(u_j\) and \(\Span\{u_1,\dots,u_{j-1}\}\) (\(2\le j\le n\)), then
\[
\boxed{\;
\det R_A=\prod_{j=2}^{n}\sin^2\theta_j,\qquad
\log\det G_A=\sum_{i=1}^{n}\log H(x_i)+2\sum_{j=2}^{n}\log(\sin\theta_j)\; }.
\]
In particular, when \(n=2\),
\[
\det R_A=1-|\langle u_1,u_2\rangle|^2=\sin^2\theta_2.
\]
\end{proposition}
\begin{proof}
From \(v(x_i)=\|v(x_i)\|u_i=\sqrt{H(x_i)}\,u_i\) we have \(G_A=DR_AD\), where \(D=\mathrm{diag}(\sqrt{H(x_1)},\dots,\sqrt{H(x_n)})\).
Taking determinants yields the first formula.
\par
Performing Gram--Schmidt on \(\{u_i\}\), let \(d_j\) be the orthogonal distance from \(u_j\) to \(\Span\{u_1,\dots,u_{j-1}\}\); then \(\det R_A=\prod_{j=1}^{n}d_j^2\) (Gram determinant equals the product of squared diagonal elements after orthogonalization), with \(d_1=1\) and \(d_j=\sin\theta_j\) (\(j\ge 2\)).
Substituting yields the displayed factorization; the two-point case is direct computation. This completes the proof.
\end{proof}

\begin{corollary}[Angle concentration in near-maximal anchoring: determinant deficit controls proportion of "non-orthogonality events"]\label{cor:xi-basepoint-scan-anchor-angle-concentration}
Following the notation of Proposition~\ref{prop:xi-basepoint-scan-anchor-principal-angles}, if
\[
\det G_A\ \ge\ e^{-\varepsilon}\prod_{i=1}^{n}H(x_i)\qquad(\varepsilon\ge 0),
\]
equivalently \(\det R_A\ge e^{-\varepsilon}\), then for any \(\eta\in(0,1)\), at least \((1-\eta)(n-1)\) indices \(j\in\{2,\dots,n\}\) satisfy
\[
\boxed{\;
\sin^2\theta_j\ \ge\ \exp\!\Bigl(-\frac{\varepsilon}{\eta(n-1)}\Bigr)\; }.
\]
\end{corollary}
\begin{proof}
By Proposition~\ref{prop:xi-basepoint-scan-anchor-principal-angles}, we have \(\prod_{j=2}^{n}\sin^2\theta_j=\det R_A\ge e^{-\varepsilon}\); taking logarithms gives
\[
\sum_{j=2}^{n}-\log(\sin^2\theta_j)\le \varepsilon.
\]
If more than \(\eta(n-1)\) indices \(j\) satisfy \(-\log(\sin^2\theta_j)>\varepsilon/(\eta(n-1))\), then the left-hand side strictly exceeds \(\varepsilon\), a contradiction.
Hence except for at most \(\eta(n-1)\) exceptional indices, the rest satisfy the displayed lower bound. This completes the proof.
\end{proof}

\begin{proposition}[\texorpdfstring{$\mathrm{PSL}(2,\RR)$}{PSL(2,R)} covariance: strict Möbius transformation law for scan kernel and anchoring determinant]\label{prop:xi-basepoint-scan-psl2r-covariance}
Under the setup of Theorem~\ref{thm:xi-basepoint-scan-finite-rank-rkhs}, take a Möbius transformation
\[
m(z)=\frac{az+b}{cz+d}\qquad(a,b,c,d\in\RR,\ ad-bc>0),
\]
which maps the upper half-plane orientation-preservingly to itself.
Let \(x_i':=m(x_i)\) and \(p_k':=m(p_k)\).
Define the transformed weights
\[
w_k':=w_k\cdot\frac{(ad-bc)^2}{|cp_k+d|^2}.
\]
Generate the transformed feature map \(v'\), kernel \(K'\) and Gram matrix \(G_{A'}'\) from \(\{(p_k',w_k')\}\).
Then for all \(x,y\in\RR\) we have the strict covariance relation
\[
\boxed{\;
K'(m(x),m(y))=(cx+d)(cy+d)\,K(x,y)\; }.
\]
Consequently, for \(A=\{x_1,\dots,x_n\}\) we have
\[
\boxed{\;
\det G_{A'}'=\Bigl(\prod_{i=1}^{n}(cx_i+d)^2\Bigr)\det G_A\; }.
\]
Moreover, the normalized correlation matrix \(R_A\) (Proposition~\ref{prop:xi-basepoint-scan-anchor-principal-angles}) is invariant under this transformation, so the "pure angular information" defined by \(\det R_A\) and the principal angle sequence \(\{\theta_j\}\) is a \(\mathrm{PSL}(2,\RR)\)-invariant.
\end{proposition}
\begin{proof}
The identity
\[
\frac{1}{m(x)-m(p)}=\frac{(cx+d)(cp+d)}{ad-bc}\cdot\frac{1}{x-p}
\]
holds for all \(x\neq p\).
Taking \(p=p_k\) and multiplying by \(\sqrt{w_k'}\) gives
\[
\frac{\sqrt{w_k'}}{m(x)-m(p_k)}
=(cx+d)\cdot \frac{cp_k+d}{|cp_k+d|}\cdot \frac{\sqrt{w_k}}{x-p_k}.
\]
Thus there exists a diagonal unitary matrix \(U=\mathrm{diag}\bigl((cp_k+d)/|cp_k+d|\bigr)\) such that for all \(x\in\RR\),
\(
v'(m(x))=(cx+d)\,U\,v(x)
\).
Taking inner products and using unitarity of \(U\) yields the kernel covariance formula; the Gram determinant scaling follows immediately from \(G_{A'}'=D G_A D\) (\(D=\mathrm{diag}(cx_i+d)\)).
Finally, the correlation matrix is invariant under common left-right scaling, so principal angles and \(\det R_A\) are invariant. This completes the proof.
\end{proof}

\begin{proposition}[Stationary-point equations for full-rank anchoring: explicit system of real-charge--complex-pole equilibrium]\label{prop:xi-basepoint-scan-full-rank-anchor-critical-equations}
In the full-rank case \(n=\kappa\) of Proposition~\ref{prop:xi-basepoint-scan-anchor-det-cauchy-vandermonde}, view \(\det G_A\) as a function of distinct real variables \((x_1,\dots,x_\kappa)\).
If \(A\) is a critical point in an open domain (partial derivative with respect to each \(x_i\) is zero), then each \(i=1,\dots,\kappa\) satisfies the equilibrium equation
\[
\boxed{\;
\sum_{\substack{j=1\\ j\ne i}}^{\kappa}\frac{1}{x_i-x_j}
=\sum_{k=1}^{\kappa}\Re\frac{1}{x_i-p_k}
=\sum_{k=1}^{\kappa}\frac{x_i-\gamma_k}{(x_i-\gamma_k)^2+a_k^2}\; }.
\]
\end{proposition}
\begin{proof}
By the closed form of Proposition~\ref{prop:xi-basepoint-scan-anchor-det-cauchy-vandermonde},
\[
\log\det G_A=\mathrm{const}+2\log|\Delta(X)|-2\sum_{i,k}\log|x_i-p_k|,
\]
where the constant term is independent of \(\{p_k,w_k\}\).
Differentiating with respect to \(x_i\) gives
\(
\partial_{x_i}\log|\Delta(X)|=\sum_{j\ne i}\frac{1}{x_i-x_j}
\)
and
\(
\partial_{x_i}\log|x_i-p_k|=\Re\frac{1}{x_i-p_k}
\).
Setting \(\partial_{x_i}\log\det G_A=0\) yields the first equality; the last formula is the real-part expansion. This completes the proof.
\end{proof}

% 2026-02-14 13:36:12 (Asia/Singapore)
\begin{proposition}[Heine--Stieltjes-type differential equation for critical anchoring configurations: existence-uniqueness of Van Vleck polynomial]\label{prop:xi-basepoint-scan-heine-stieltjes-ode}
In the full-rank case \(n=\kappa\) of Proposition~\ref{prop:xi-basepoint-scan-full-rank-anchor-critical-equations}, suppose the anchor points \(X=\{x_1,\dots,x_\kappa\}\subset\RR\) are mutually distinct, and let the anchor polynomial
\(
y(z):=\prod_{j=1}^{\kappa}(z-x_j)
\)
be monic.
Let the pole polynomial be
\[
\mathcal{A}(z):=\prod_{k=1}^{\kappa}(z-p_k)(z-\overline{p_k})\in\RR[z].
\]
Then the following are equivalent:
\begin{enumerate}
\item The anchor configuration \(X\) is a critical configuration of \(\det G_A\) (equivalently satisfies the equilibrium equations in Proposition~\ref{prop:xi-basepoint-scan-full-rank-anchor-critical-equations}).
\item There exists a polynomial \(B\in\RR[z]\) (\(\deg B\le 2\kappa-2\)) such that the identity
\[
\boxed{\;
\mathcal{A}(z)\,y''(z)-\mathcal{A}'(z)\,y'(z)=B(z)\,y(z)\;}
\]
holds identically on \(\CC\).
\end{enumerate}
Moreover, when this identity holds, the polynomial \(B\) is uniquely determined by \(y\) (called the corresponding Van Vleck polynomial).
\end{proposition}
\begin{proof}
Denote
\(
N(z):=\mathcal{A}(z)y''(z)-\mathcal{A}'(z)y'(z)
\).
\par
\emph{(1)\(\Rightarrow\)(2).}
For any root \(x_i\) (so \(y(x_i)=0\) and \(y'(x_i)\ne 0\)) we have the identity
\[
\frac{y''(x_i)}{y'(x_i)}=2\sum_{j\ne i}\frac{1}{x_i-x_j}.
\]
On the other hand,
\[
\frac{\mathcal{A}'(z)}{\mathcal{A}(z)}
=\sum_{k=1}^{\kappa}\left(\frac{1}{z-p_k}+\frac{1}{z-\overline{p_k}}\right),
\]
so when \(z=x_i\in\RR\)
\(
\frac{\mathcal{A}'(x_i)}{\mathcal{A}(x_i)}
=2\sum_{k=1}^{\kappa}\Re\frac{1}{x_i-p_k}
\).
Combining this with the critical equilibrium equation
\(
\sum_{j\ne i}\frac{1}{x_i-x_j}=\sum_{k}\Re\frac{1}{x_i-p_k}
\)
yields
\(
y''(x_i)/y'(x_i)=\mathcal{A}'(x_i)/\mathcal{A}(x_i)
\), equivalently
\[
N(x_i)=\mathcal{A}(x_i)y''(x_i)-\mathcal{A}'(x_i)y'(x_i)=0.
\]
Since \(\{x_i\}\) are mutually distinct and \(y\) is a monic polynomial of degree \(\kappa\), we have \(y\mid N\).
Let \(B:=N/y\); then \(B\in\RR[z]\), and
\(
\deg B\le \deg N-\deg y\le (2\kappa+\kappa-2)-\kappa=2\kappa-2
\).
\par
\emph{(2)\(\Rightarrow\)(1).}
Suppose \(\mathcal{A}y''-\mathcal{A}'y'=By\) holds.
Substituting any root \(x_i\) gives
\(
\mathcal{A}(x_i)y''(x_i)-\mathcal{A}'(x_i)y'(x_i)=0
\).
Since
\(
\mathcal{A}(x_i)=\prod_{k}|x_i-p_k|^2>0
\),
we can divide to get
\(
y''(x_i)/y'(x_i)=\mathcal{A}'(x_i)/\mathcal{A}(x_i)
\).
Substituting the partial fraction expansion of \(\mathcal{A}'/\mathcal{A}\) and using
\(
y''(x_i)/y'(x_i)=2\sum_{j\ne i}\frac{1}{x_i-x_j}
\),
we obtain the critical equilibrium equation
\(
\sum_{j\ne i}\frac{1}{x_i-x_j}=\sum_k\Re\frac{1}{x_i-p_k}
\).
By Proposition~\ref{prop:xi-basepoint-scan-full-rank-anchor-critical-equations}, this is equivalent to \(X\) being a critical configuration of \(\det G_A\).
\par
Finally, if \(\mathcal{A}y''-\mathcal{A}'y'=B_1y=B_2y\), then \((B_1-B_2)y\equiv 0\) forces \(B_1=B_2\).
Hence \(B\) is unique. This completes the proof.
\end{proof}

\begin{proposition}[Leading-coefficient rigidity and sub-leading closed form of Van Vleck polynomial]\label{prop:xi-basepoint-scan-van-vleck-leading-coeffs}
Under the setup of Proposition~\ref{prop:xi-basepoint-scan-heine-stieltjes-ode}, suppose \(y\) is monic and \(\deg y=\kappa\), and let
\[
B(z)=b_0 z^{2\kappa-2}+b_1 z^{2\kappa-3}+O(z^{2\kappa-4}),
\quad
\mathcal{A}(z)=z^{2\kappa}+\mathcal{A}_1 z^{2\kappa-1}+O(z^{2\kappa-2}),
\quad
y(z)=z^{\kappa}+y_1 z^{\kappa-1}+O(z^{\kappa-2}).
\]
Then necessarily
\[
\boxed{\;b_0=-\kappa(\kappa+1),\qquad b_1=-\kappa^{2}\mathcal{A}_1+2y_1.\;}
\]
Here
\(
\mathcal{A}_1=-\sum_{k=1}^{\kappa}(p_k+\overline{p_k})=-2\sum_{k=1}^{\kappa}\gamma_k
\),
and
\(
y_1=-\sum_{j=1}^{\kappa}x_j
\).
\end{proposition}
\begin{proof}
Compare leading terms of the identity \(\mathcal{A}y''-\mathcal{A}'y'=By\) at infinity.
From
\[
y'(z)=\kappa z^{\kappa-1}+(\kappa-1)y_1 z^{\kappa-2}+O(z^{\kappa-3}),
\qquad
y''(z)=\kappa(\kappa-1)z^{\kappa-2}+(\kappa-1)(\kappa-2)y_1 z^{\kappa-3}+O(z^{\kappa-4}),
\]
and
\[
\mathcal{A}'(z)=2\kappa z^{2\kappa-1}+(2\kappa-1)\mathcal{A}_1 z^{2\kappa-2}+O(z^{2\kappa-3}),
\]
we obtain
\[
\mathcal{A}y''-\mathcal{A}'y'
=\bigl(\kappa(\kappa-1)-2\kappa^2\bigr)z^{3\kappa-2}
\bigl(-\kappa^2\mathcal{A}_1-(\kappa-1)(\kappa+2)y_1\bigr)z^{3\kappa-3}
+O(z^{3\kappa-4}).
\]
On the other hand,
\[
By=b_0 z^{3\kappa-2}+(b_1+b_0y_1)z^{3\kappa-3}+O(z^{3\kappa-4}).
\]
Comparing term-by-term yields \(b_0=-\kappa(\kappa+1)\) and \(b_1=-\kappa^{2}\mathcal{A}_1+2y_1\). This completes the proof.
\end{proof}

\begin{proposition}[Zeroth- and first-order residue moment laws for Van Vleck ratio: strict linear constraints on defect current]\label{prop:xi-basepoint-scan-van-vleck-residue-moment01}
Under the setup of Proposition~\ref{prop:xi-basepoint-scan-heine-stieltjes-ode}, let
\(
f(z):=\frac{B(z)}{\mathcal{A}(z)}
\).
Then \(f\) is a rational function with simple poles only at \(\{p_k,\overline{p_k}\}\), and by Proposition~\ref{prop:xi-basepoint-scan-van-vleck-leading-coeffs} we have
\(
f(z)=-\kappa(\kappa+1)z^{-2}+O(z^{-3})
\).
Let
\[
r_k:=\operatorname*{Res}_{z=p_k} f(z)=\frac{B(p_k)}{\mathcal{A}'(p_k)}.
\]
Then necessarily
\[
\boxed{\;
\sum_{k=1}^{\kappa}\Re r_k=0,\qquad
\sum_{k=1}^{\kappa}\Re(p_k r_k)=-\frac{\kappa(\kappa+1)}{2}\; }.
\]
If \(\mathcal{A}y''-\mathcal{A}'y'=By\) holds, then further at each \(p_k\) we have
\[
r_k=-\frac{y'(p_k)}{y(p_k)},
\]
hence equivalently
\[
\boxed{\;
\sum_{k=1}^{\kappa}\Re\frac{y'(p_k)}{y(p_k)}=0,\qquad
\sum_{k=1}^{\kappa}\Re\!\left(p_k\frac{y'(p_k)}{y(p_k)}\right)=\frac{\kappa(\kappa+1)}{2}\; }.
\]
\end{proposition}
\begin{proof}
Since \(\deg B=\deg\mathcal{A}-2\), \(f\) has no polynomial part at infinity.
Moreover, since \(B,\mathcal{A}\in\RR[z]\), the poles of \(f\) come in conjugate pairs, and its partial fraction expansion can be written as
\[
f(z)=\sum_{k=1}^{\kappa}\frac{r_k}{z-p_k}+\sum_{k=1}^{\kappa}\frac{\overline{r_k}}{z-\overline{p_k}}.
\]
Expanding at infinity gives
\[
f(z)=\frac{\sum_k(r_k+\overline{r_k})}{z}
\;+\;\frac{\sum_k(p_k r_k+\overline{p_k}\,\overline{r_k})}{z^2}
\;+\;O(z^{-3}).
\]
By Proposition~\ref{prop:xi-basepoint-scan-van-vleck-leading-coeffs},
\(
f(z)=-\kappa(\kappa+1)z^{-2}+O(z^{-3})
\), the \(z^{-1}\) coefficient is zero, so \(\sum_k\Re r_k=0\).
The \(z^{-2}\) coefficient is \(-\kappa(\kappa+1)\), so
\(
2\sum_k\Re(p_k r_k)=-\kappa(\kappa+1)
\).
\par
Finally, if \(\mathcal{A}y''-\mathcal{A}'y'=By\) holds, then at \(z=p_k\) where \(\mathcal{A}(p_k)=0\) we have
\[
-\mathcal{A}'(p_k)y'(p_k)=B(p_k)y(p_k),
\]
hence \(r_k=B(p_k)/\mathcal{A}'(p_k)=-y'(p_k)/y(p_k)\). This completes the proof.
\end{proof}
\begin{proposition}[Second-order residue moment law and explicit "center-of-mass term": closed identity for pole-square weighting]\label{prop:xi-basepoint-scan-van-vleck-residue-moment2}
Following the notation of Proposition~\ref{prop:xi-basepoint-scan-van-vleck-residue-moment01}, we have the strict identity
\[
\boxed{\;
\sum_{k=1}^{\kappa}\Re(p_k^{2} r_k)=\frac{1}{2}\bigl(b_1-b_0\mathcal{A}_1\bigr)\; }.
\]
In the case \(r_k=-y'(p_k)/y(p_k)\), combining with Proposition~\ref{prop:xi-basepoint-scan-van-vleck-leading-coeffs}'s
\(
b_0=-\kappa(\kappa+1)
\),
\(
b_1=-\kappa^{2}\mathcal{A}_1+2y_1
\),
and using \(\mathcal{A}_1=-2\sum_{k}\gamma_k\), \(y_1=-\sum_{j}x_j\), we simplify to obtain
\[
\boxed{\;
\sum_{k=1}^{\kappa}\Re\!\left(p_k^{2}\frac{y'(p_k)}{y(p_k)}\right)
=\kappa\sum_{k=1}^{\kappa}\gamma_k+\sum_{j=1}^{\kappa}x_j\; }.
\]
\end{proposition}
\begin{proof}
By the partial fraction expansion in Proposition~\ref{prop:xi-basepoint-scan-van-vleck-residue-moment01}, expanding at infinity gives the \(z^{-3}\) coefficient as
\[
\sum_{k=1}^{\kappa}(p_k^2 r_k+\overline{p_k}^2\,\overline{r_k})
=2\sum_{k=1}^{\kappa}\Re(p_k^2 r_k).
\]
On the other hand,
\[
\frac{B(z)}{\mathcal{A}(z)}
=b_0 z^{-2}+(b_1-b_0\mathcal{A}_1)z^{-3}+O(z^{-4}),
\]
comparing the \(z^{-3}\) coefficient yields the first formula.
\par
If \(r_k=-y'(p_k)/y(p_k)\), substituting this together with Proposition~\ref{prop:xi-basepoint-scan-van-vleck-leading-coeffs}'s \(b_0,b_1\) into the first formula and using \(\mathcal{A}_1=-2\sum\gamma_k\), \(y_1=-\sum x_j\) yields the second formula. This completes the proof.
\end{proof}
\begin{proposition}[Existence of global maximum for full-rank anchoring determinant and necessity of critical points]\label{prop:xi-basepoint-scan-anchor-det-global-maximum}
In the full-rank case \(n=\kappa\) of Proposition~\ref{prop:xi-basepoint-scan-anchor-det-cauchy-vandermonde}, view \(\det G_A\) as a continuous function on the real configuration space
\[
\Omega:=\{(x_1,\dots,x_\kappa)\in\RR^\kappa:\ x_i\neq x_j\}
\]
(permutation-symmetric).
Then \(\det G_A\) attains a global maximum; for any maximum point \(X^\star\in\Omega\), it must be a critical configuration, hence satisfying the equilibrium equations in Proposition~\ref{prop:xi-basepoint-scan-full-rank-anchor-critical-equations}.
Consequently, by Proposition~\ref{prop:xi-basepoint-scan-heine-stieltjes-ode}, there exists a corresponding monic anchor polynomial \(y^\star\) and unique Van Vleck polynomial \(B^\star\) such that
\[
\mathcal{A}(z)\,(y^\star)''(z)-\mathcal{A}'(z)\,(y^\star)'(z)=B^\star(z)\,y^\star(z)
\]
holds.
\end{proposition}
\begin{proof}
By Theorem~\ref{thm:xi-basepoint-scan-finite-rank-rkhs}, for any distinct anchor points \(X\in\Omega\), the Gram matrix \(G_A\) is positive definite, so \(\det G_A>0\).
Let \(M:=\sup_{X\in\Omega}\det G_A\in(0,\infty)\), and take a maximizing sequence \(X^{(m)}\in\Omega\) such that \(\det G_{A^{(m)}}\to M\).
\par
First we rule out escape to infinity.
By Hadamard's inequality and \(G_{ii}=H(x_i)\) (Theorem~\ref{thm:xi-basepoint-scan-finite-rank-rkhs}), we have
\[
\det G_A\le \prod_{i=1}^{\kappa}H(x_i).
\]
Since \(H(x)=\sum_k w_k/|x-p_k|^2=O(x^{-2})\) (\(|x|\to\infty\)), if \(\max_i|x_i^{(m)}|\to\infty\), then the right-hand side tends to \(0\), contradicting \(\det G_{A^{(m)}}\to M>0\).
Thus \(\{X^{(m)}\}\) is bounded, so there exists a convergent subsequence (still denoted \(X^{(m)}\)) such that \(X^{(m)}\to X^\star\in\RR^\kappa\).
\par
Next we rule out collision boundary.
By the full-rank closed form of Proposition~\ref{prop:xi-basepoint-scan-anchor-det-cauchy-vandermonde},
\(
\det G_A\propto |\Delta(X)|^2
\);
when a collision \(x_i\to x_j\) occurs, \(\Delta(X)\to 0\), so \(\det G_A\to 0\).
Thus the limit point \(X^\star\) of the maximizing sequence cannot lie on the collision boundary; it must satisfy \(X^\star\in\Omega\).
By continuity, \(\det G_{A^\star}=M\), i.e., the global maximum is attained.
\par
Since \(\Omega\) is an open set, the maximum point \(X^\star\) is an interior point, so \(\nabla \det G_A(X^\star)=0\), i.e., it is a critical configuration.
By Proposition~\ref{prop:xi-basepoint-scan-full-rank-anchor-critical-equations} the equilibrium equations hold; finally by Proposition~\ref{prop:xi-basepoint-scan-heine-stieltjes-ode} we obtain the corresponding differential equation and existence-uniqueness of \(B^\star\). This completes the proof.
\end{proof}
\begin{proposition}[Analytic rigidity of critical anchoring configurations: real-analytic dependence on pole parameters and local isolation]\label{prop:xi-basepoint-scan-critical-configuration-analytic-rigidity}
Under the full-rank case of Proposition~\ref{prop:xi-basepoint-scan-full-rank-anchor-critical-equations}, suppose the pole multiset \(\{p_k\}\subset\CC_{-}\) is distinct.
Consider the map \(\Phi=(\Phi_1,\dots,\Phi_\kappa)\):
\[
\Phi_i(x_1,\dots,x_\kappa;\,p_1,\dots,p_\kappa)
:=\sum_{j\ne i}\frac{1}{x_i-x_j}-\sum_{k=1}^{\kappa}\Re\frac{1}{x_i-p_k},
\qquad i=1,\dots,\kappa.
\]
If there exists a critical configuration \(X^\star=(x_1^\star,\dots,x_\kappa^\star)\) satisfying \(\Phi(X^\star;p)=0\), and the Jacobian at this point
\(
J:=\bigl(\partial \Phi_i/\partial x_j\bigr)_{1\le i,j\le \kappa}
\)
is invertible, then in a sufficiently small neighborhood of \(p\) there exists a unique real-analytic branch \(X(p)\) such that \(\Phi(X(p);p)=0\) and \(X(p^\star)=X^\star\).
Hence under non-degeneracy conditions, the critical anchoring configuration exhibits real-analytic rigidity with respect to pole parameters and is locally isolated in that neighborhood.
\end{proposition}
\begin{proof}
\(\Phi\) is real-analytic in \((x,p)\) on the domain of distinct points.
The invertible Jacobian hypothesis allows application of the implicit function theorem, hence there exists a unique real-analytic solution branch \(X(p)\) yielding local isolation. This completes the proof.
\end{proof}
% (User input window)
% 2026-02-14 16:12:45 (Asia/Singapore)

\paragraph{Completeness relation and Gram inverse closed form for finite-rank anchoring kernel: full reconstruction of continuous kernel from discrete anchor points}
Following the notation of Theorem~\ref{thm:xi-basepoint-scan-finite-rank-rkhs}: poles \(\{p_k\}_{k=1}^{\kappa}\subset\CC_{-}\), weights \(\{w_k\}_{k=1}^{\kappa}\subset\RR_{>0}\), feature map \(v\) and kernel \(K\).

\begin{proposition}[Kernel completeness relation for full-rank anchor sets: full reconstruction of continuous kernel from discrete anchors]\label{prop:xi-basepoint-scan-kernel-completeness-full-rank}
For any \(x\in\RR\) define
\[
v(x)_k:=\frac{\sqrt{w_k}}{x-p_k}\in\CC,\qquad
K(x,y):=\langle v(x),v(y)\rangle_{\CC^\kappa}
=\sum_{k=1}^{\kappa}\frac{w_k}{(x-p_k)(y-\overline{p_k})}\qquad(x,y\in\RR).
\]
Take distinct real anchor points \(A=\{a_1,\dots,a_\kappa\}\subset\RR\), and denote the Gram matrix
\[
G_A:=(K(a_i,a_j))_{1\le i,j\le \kappa}.
\]
Then \(G_A\) is invertible, and for all \(x,y\in\RR\) the completeness identity holds:
\[
\boxed{\;
K(x,y)=K(x,A)\,G_A^{-1}\,K(A,y)\; },
\]
where
\(
K(x,A)=(K(x,a_1),\dots,K(x,a_\kappa))
\),
\(
K(A,y)=(K(a_1,y),\dots,K(a_\kappa,y))^{\mathsf T}
\).
\end{proposition}
\begin{proof}
Let the matrix \(V\in\CC^{\kappa\times\kappa}\) be
\(
V_{ik}:=v(a_i)_k=\sqrt{w_k}/(a_i-p_k)
\).
Then \(G_A=VV^\ast\).
By the Cauchy non-degeneracy in the proof of Theorem~\ref{thm:xi-basepoint-scan-finite-rank-rkhs}, \(V\) is invertible, so
\[
I_\kappa=V^\ast G_A^{-1}V.
\]
Multiplying both sides on the left by \(v(x)^\ast\) and on the right by \(v(y)\) gives
\[
K(x,y)
=\langle v(x),v(y)\rangle
=\langle v(x),V^\ast G_A^{-1}Vv(y)\rangle
=K(x,A)\,G_A^{-1}\,K(A,y),
\]
which is the displayed completeness relation. This completes the proof.
\end{proof}

\begin{proposition}[Full closed form of full-rank Gram inverse: \texorpdfstring{$G_A^{-1}$}{G\_A^{-1}} given explicitly via Cauchy inverse]\label{prop:xi-basepoint-scan-gram-inverse-closed-form}
Under the full-rank setup of Proposition~\ref{prop:xi-basepoint-scan-kernel-completeness-full-rank}, let
\[
\Pi(z):=\prod_{k=1}^{\kappa}(z-p_k),\qquad
y(z):=\prod_{j=1}^{\kappa}(z-a_j).
\]
Write \(V_{ik}=\sqrt{w_k}/(a_i-p_k)\), so \(V=C\,\mathrm{diag}(\sqrt w)\), where \(C_{ik}=(a_i-p_k)^{-1}\).
By the Cauchy inverse closed form of Conclusion~\ref{con:xi-terminal-cauchy-lagrange-rational-interpolation-cauchy-inverse}, we have
\[
(V^{-1})_{k i}
=\frac{1}{\sqrt{w_k}}\,(C^{-1})_{k i}
=\frac{1}{\sqrt{w_k}}\cdot
\frac{y(p_k)\,\Pi(a_i)}{(p_k-a_i)\,\Pi'(p_k)\,y'(a_i)}.
\]
Consequently
\[
G_A^{-1}=(V^{-1})^\ast V^{-1},
\]
and its entries have the explicit expression: for \(1\le i,j\le \kappa\),
\[
\boxed{\;
(G_A^{-1})_{ij}
=
\frac{\overline{\Pi(a_i)}\,\Pi(a_j)}{y'(a_i)\,y'(a_j)}
\sum_{k=1}^{\kappa}
\frac{|y(p_k)|^2}{w_k\,|\Pi'(p_k)|^2}\cdot
\frac{1}{(a_i-\overline{p_k})(a_j-p_k)}\; }.
\]
\end{proposition}
\begin{proof}
Since \(G_A=VV^\ast\) and \(V\) is invertible, we have \(G_A^{-1}=(V^{-1})^\ast V^{-1}\).
Moreover, \(V=C\,\mathrm{diag}(\sqrt w)\) gives \(V^{-1}=\mathrm{diag}(w^{-1/2})\,C^{-1}\); substituting the \(C^{-1}\) closed form from Conclusion~\ref{con:xi-terminal-cauchy-lagrange-rational-interpolation-cauchy-inverse} yields \((V^{-1})_{ki}\).
Expanding
\(
(G_A^{-1})_{ij}=\sum_{k}\overline{(V^{-1})_{ki}}(V^{-1})_{kj}
\)
term-by-term and simplifying gives the displayed expression. This completes the proof.
\end{proof}

\begin{proposition}[Determinant increment and orthogonal gap: geometric meaning of Schur complement]\label{prop:xi-basepoint-scan-gram-determinant-schur-increment}
For any anchor set \(A=\{a_1,\dots,a_n\}\subset\RR\) (\(n<\kappa\)) assuming \(G_A:=(K(a_i,a_j))_{i,j}\) is invertible, define
\[
k_A(x):=(K(x,a_1),\dots,K(x,a_n))^{\mathsf T},\qquad
H(x):=K(x,x).
\]
Then for any \(x\in\RR\setminus A\) we have the strict identity
\[
\boxed{\;
\det G_{A\cup\{x\}}=\det G_A\cdot\Bigl(H(x)-k_A(x)^\ast G_A^{-1}k_A(x)\Bigr)\; } .
\]
Moreover, the quantity inside the parentheses is non-negative and satisfies the orthogonal-gap representation
\[
H(x)-k_A(x)^\ast G_A^{-1}k_A(x)
=\bigl\|v(x)-\mathrm{proj}_{\Span\{v(a_i)\}}v(x)\bigr\|_{\CC^\kappa}^{2}.
\]
\end{proposition}
\begin{proof}
Write in block-matrix form
\[
G_{A\cup\{x\}}=
\begin{pmatrix}
G_A & k_A(x)\\
k_A(x)^\ast & H(x)
\end{pmatrix}.
\]
Applying the Schur complement formula to the block determinant yields the first formula.
The second formula follows from \(G_A=V_AV_A^\ast\) (\(V_A\) is the \(n\times\kappa\) matrix formed by \(v(a_i)\)) and the least-squares projection identity:
\(
k_A(x)=V_A v(x)
\), so the Schur complement equals the squared orthogonal-gap norm of \(v(x)\) from \(\Span\{v(a_i)\}\). This completes the proof.
\end{proof}

\begin{corollary}[Leverage factorization of Gram determinant: multiplicative factorization and global upper bound under finite-rank kernel]\label{cor:xi-basepoint-scan-gram-determinant-leverage-product}
For any ordered anchor sequence \((a_1,\dots,a_n)\) (distinct, with each \(G_{A_m}\) invertible), denote the incremental anchor sets \(A_m:=\{a_1,\dots,a_m\}\) (\(A_0=\varnothing\)), and define the leverage gap
\[
\ell_{A_{m-1}}(a_m):=
H(a_m)-k_{A_{m-1}}(a_m)^\ast G_{A_{m-1}}^{-1}k_{A_{m-1}}(a_m)\ \ge 0.
\]
Then the strict multiplicative factorization holds:
\[
\boxed{\;
\det G_{A_n}=\prod_{m=1}^{n}\ell_{A_{m-1}}(a_m)\; }.
\]
Consequently, for any \(n\)-point anchor set \(A\) (ignoring order) we have the upper bound
\[
\det G_A\le \prod_{m=1}^{n}\Bigl(\sup_{x\in\RR}\ell_{A_{m-1}}(x)\Bigr),
\]
with equality if and only if there exists a permutation such that each step \(a_m\) achieves the corresponding supremum.
\end{corollary}
\begin{proof}
Applying Proposition~\ref{prop:xi-basepoint-scan-gram-determinant-schur-increment} for \(m=1,\dots,n\) successively gives
\[
\det G_{A_m}=\det G_{A_{m-1}}\cdot \ell_{A_{m-1}}(a_m),
\]
multiplying yields the factorization.
The upper bound follows from \(\ell_{A_{m-1}}(a_m)\le \sup_x\ell_{A_{m-1}}(x)\) multiplied term-by-term; the equality condition follows from the necessity-sufficiency of achieving equality at each step. This completes the proof.
\end{proof}


\endinput








